\chapter{Wave Equation}

\section*{Introduction}

Every time you hear music, see a rainbow, feel an earthquake, or use Wi-Fi, you are experiencing the wave equation in action. It says that any disturbance in a medium propagates at a speed determined by the medium's properties---tension and density for strings, pressure and compressibility for air, permittivity and permeability for light. The same equation connects the vibrations of a violin string to the oscillations of distant quasars.

\begin{enumerate}
\item \textbf{Wave Equation (1D):} The displacement of a wave satisfies a second-order PDE.

\item \textbf{3D Wave Equation:} Extends to three dimensions.

\end{enumerate}
\section*{Fundamental Equation Used}

\[
\frac{\partial^2 u}{\partial t^2} = c^2 \frac{\partial^2 u}{\partial x^2}
\]


\section*{Assumptions}

\textbf{Assumptions:} The medium is linear, homogeneous, and non-dispersive (wave speed $c$ is constant). Small-amplitude disturbances (linearization). No damping or energy loss.


\textbf{Limitations:} Does not account for dispersion (frequency-dependent speed, needed for waves in optical fibers or deep water). Does not include nonlinear effects (shock waves, solitons). Does not account for absorption or scattering losses.


\section*{Mathematical Derivation}

\textbf{Step 1: The Physical Setup---A Vibrating String.}

Consider a perfectly flexible string under uniform tension $T$, with linear mass density $\mu$ (mass per unit length). The string is displaced from equilibrium by a small amount $u(x,t)$ in the transverse direction. We assume small displacements ($|u| \ll L$), so the tension remains approximately constant and the motion is purely transverse.

\textbf{Step 2: Forces on a Small Element.}

Isolate an infinitesimal element of the string between positions $x$ and $x + \Delta x$. The tension $T$ acts tangentially at each end of the element. At position $x$, the tension vector makes angle $\alpha$ with the horizontal; at $x + \Delta x$, it makes angle $\beta$:

\begin{equation}
\text{Net vertical force} = T\sin\beta - T\sin\alpha.
\end{equation}

For small angles, $\sin\alpha \approx \tan\alpha = \partial u/\partial x\big|_x$ and $\sin\beta \approx \tan\beta = \partial u/\partial x\big|_{x+\Delta x}$:

\begin{equation}
F_{\text{net}} = T\left[\frac{\partial u}{\partial x}\bigg|_{x+\Delta x} - \frac{\partial u}{\partial x}\bigg|_x\right] \approx T\,\frac{\partial^2 u}{\partial x^2}\,\Delta x.
\end{equation}

\textbf{Step 3: Newton's Second Law.}

The mass of the element is $\mu\,\Delta x$, and its transverse acceleration is $\partial^2 u/\partial t^2$. Applying $F = ma$:

\begin{equation}
\mu\,\Delta x\,\frac{\partial^2 u}{\partial t^2} = T\,\frac{\partial^2 u}{\partial x^2}\,\Delta x.
\end{equation}

Canceling $\Delta x$:

\begin{equation}
\mu\,\frac{\partial^2 u}{\partial t^2} = T\,\frac{\partial^2 u}{\partial x^2}.
\end{equation}

\textbf{Step 4: The Standard Form.}

Dividing by $\mu$ and defining $c^2 = T/\mu$:

\begin{equation}
\boxed{\frac{\partial^2 u}{\partial t^2} = c^2\,\frac{\partial^2 u}{\partial x^2}.}
\end{equation}

This is the \textbf{one-dimensional wave equation}. The wave speed is $c = \sqrt{T/\mu}$, determined by the competition between the restoring force (tension $T$) and the inertia (mass density $\mu$).

\textbf{Step 5: The d'Alembert Solution.}

Introduce the change of variables $\xi = x - ct$ and $\eta = x + ct$. By the chain rule:

\begin{align}
\frac{\partial u}{\partial x} &= \frac{\partial u}{\partial \xi} + \frac{\partial u}{\partial \eta}, \qquad \frac{\partial^2 u}{\partial x^2} = \frac{\partial^2 u}{\partial \xi^2} + 2\frac{\partial^2 u}{\partial \xi\,\partial \eta} + \frac{\partial^2 u}{\partial \eta^2}, \\
\frac{\partial u}{\partial t} &= -c\frac{\partial u}{\partial \xi} + c\frac{\partial u}{\partial \eta}, \qquad \frac{\partial^2 u}{\partial t^2} = c^2\frac{\partial^2 u}{\partial \xi^2} - 2c^2\frac{\partial^2 u}{\partial \xi\,\partial \eta} + c^2\frac{\partial^2 u}{\partial \eta^2}.
\end{align}

Substituting into the wave equation:

\begin{equation}
c^2\left(\frac{\partial^2 u}{\partial \xi^2} - 2\frac{\partial^2 u}{\partial \xi\,\partial \eta} + \frac{\partial^2 u}{\partial \eta^2}\right) = c^2\left(\frac{\partial^2 u}{\partial \xi^2} + 2\frac{\partial^2 u}{\partial \xi\,\partial \eta} + \frac{\partial^2 u}{\partial \eta^2}\right).
\end{equation}

Simplifying: $-4c^2\,\partial^2 u/(\partial\xi\,\partial\eta) = 0$, which gives:

\begin{equation}
\frac{\partial^2 u}{\partial \xi\,\partial \eta} = 0.
\end{equation}

Integrating: $u(\xi, \eta) = f(\xi) + g(\eta)$, or equivalently:

\begin{equation}
\boxed{u(x,t) = f(x - ct) + g(x + ct).}
\end{equation}

This is d'Alembert's solution: \emph{any} function of $(x - ct)$ is a right-traveling wave, and any function of $(x + ct)$ is a left-traveling wave. The wave maintains its shape while propagating at speed $c$.

\textbf{Step 6: Standing Waves and Normal Modes.}

For a string fixed at both ends ($u(0,t) = u(L,t) = 0$), the boundary conditions restrict the solutions to standing waves:

\begin{equation}
u_n(x,t) = A_n\sin\!\left(\frac{n\pi x}{L}\right)\cos(\omega_n t + \phi_n), \qquad \omega_n = \frac{n\pi c}{L}, \qquad n = 1, 2, 3, \ldots
\end{equation}

Each normal mode has a discrete frequency $\omega_n = n\omega_1$, forming the harmonic series---the mathematical basis of musical harmony.

\section*{Characteristics of the Equation}

The d'Alembert solution: $u(x,t) = f(x - ct) + g(x + ct)$---any function of $(x - ct)$ and $(x + ct)$ is a solution. Standing waves: $u = A \sin(kx)\cos(\omega t)$ with $\omega = ck$. Energy is equally divided between kinetic and potential.
\section*{Solution Methods}

d'Alembert's traveling wave solution, separation of variables (standing waves), Fourier analysis (decompose into normal modes), numerical finite-difference methods (Yee algorithm), Green's functions for initial value problems.
\section*{Verifications}

Standing wave resonance on strings (verify $f_n = nc/2L$). Speed of sound measurements in tubes. Interference experiments (Young's slits). Doppler shift measurements confirm wave propagation speed.
\section*{Applications}

Acoustics (room design, musical instruments), seismology (earthquake wave analysis), telecommunications (radio, microwave, fiber optics), medical ultrasound, structural dynamics (building vibrations), oceanography (ocean waves).
\section*{What Richard Feynman Would Have Asked}

\subsection*{1. The Plain-English Translation}
\textit{``What is the wave equation really telling us about the world?''}

The Core Meaning: If you disturb a medium, the disturbance propagates outward at a speed set by the medium. The equation says that the acceleration of any point in the medium is proportional to how much it differs from the average of its neighbors. Points that are ``above'' their neighbors accelerate downward; points that are ``below'' accelerate upward. This competition between inertia and restoring force produces waves.

The Metaphor: Imagine a line of people holding hands. If one person is pushed, the push propagates down the line. How fast it travels depends on how tightly they hold hands (stiffness) and how heavy they are (inertia). The wave equation quantifies this: $c = \sqrt{\text{stiffness}/\text{inertia}}$.

\subsection*{2. The Localized Mechanics}
\textit{``At a single point on the string, what forces are at play?''}

He would press you on the physics behind the equation. At each point, the tension pulls toward the average position of the neighbors (the $\partial^2 u/\partial x^2$ term). The inertia resists acceleration (the $\partial^2 u/\partial t^2$ term). The balance between these two gives the wave equation. If the curvature is positive (the point is below its neighbors), it accelerates upward, and vice versa.

The Smoothness Check: He would ask why the equation requires second derivatives in both space and time. The second time derivative is Newton's second law ($F = ma$). The second space derivative measures curvature, which determines the restoring force (tension times curvature). First-order derivatives would give diffusion (damping), not waves.

\subsection*{3. Testing the Extreme Limits}
\textit{``What happens when we push the wave to extreme conditions?''}

The High-Frequency Limit: As frequency increases, wavelength decreases. Eventually, the wavelength approaches the spacing between atoms in the medium---the continuum approximation breaks down and the wave equation no longer applies. This is the phonon cutoff in solids.

The Large-Amplitude Limit: When amplitudes become large, the linear approximation fails. Nonlinear wave equations (Korteweg-de Vries, nonlinear Schr\"odinger) are needed. Shock waves, solitons, and harmonic generation are nonlinear phenomena.

The Vacuum Limit: In electromagnetic waves, the ``medium'' is the vacuum itself. The wave equation still holds---light propagates through empty space because the vacuum has electrical and magnetic properties ($\epsilon_0$ and $\mu_0$).

\subsection*{4. Dimensionality and Geometry}
\textit{``How does the number of dimensions change wave behavior?''}

He would point out that in 1D, waves travel without losing amplitude. In 2D (waves on a drumhead), amplitude decreases as $1/\sqrt{r}$ (cylindrical spreading). In 3D (sound in air), amplitude decreases as $1/r$ (spherical spreading). This geometric spreading is why you can hear someone nearby but not someone far away---the energy is spread over a larger sphere.

\subsection*{5. Cross-Disciplinary Connection}
\textit{``Where else does the wave equation appear?''}

The wave equation appears in: financial mathematics (option pricing near boundaries), population dynamics (migration waves), neural science (action potential propagation), social science (spread of information), and quantum mechanics (the Schr\"odinger equation is a modified wave equation with a first time derivative). Any system where a restoring force competes with inertia produces waves.
\section*{FAQ}

\textbf{Q: Is light really described by the wave equation?}\\
\textbf{A:} Yes. Maxwell's equations reduce to the wave equation for the electric and magnetic fields in vacuum, with $c = 1/\sqrt{\mu_0 \epsilon_0} \approx 3 \times 10^8$ m/s. This was Maxwell's greatest prediction.

\textbf{Q: Can waves pass through each other?}\\
\textbf{A:} Yes. The wave equation is linear, so the sum of two solutions is also a solution. Waves pass through each other without distortion (superposition principle). This is why you can hear multiple sounds simultaneously.
\section*{Important Bibliography}

\begin{enumerate}
\item d'Alembert, J. le R., ``Recherches sur la courbe qui forme une corde tendue et vibrante,'' \textit{Acad\'emie Royale des Sciences et Belles-Lettres de Berlin} \textbf{3}, 214--219 (1747).
\item Stakgold, I. and Holst, M.J., \textit{Green's Functions and Boundary Value Problems}, 3rd ed. (Wiley, 2011), Chapters 6--8.
\item Courant, R. and Hilbert, D., \textit{Methods of Mathematical Physics}, Vol. I (Wiley-Interscience, 1989), Chapters 5--6.
\item Bender, C.M. and Orszag, S.A., \textit{Advanced Mathematical Methods for Scientists and Engineers} (Springer, 1999), Chapters 5--6.
\end{enumerate}


